% page 86 du fac-simile (image p-085 du scan)
% bloc 160.0 x 242.0 mm ; marge gauche 8.0 mm
\pgc{-12.700mm}{0.000mm}{
\l{}
\l{}
\l{\cel{9}78}
\l{}
\l{\sou{burjono}. (\sou{bot.}) Sproso, rudimenta, qua kontenas jerme la stipi,}
\l{\cel{3}branchi, folii, flori, e frukti di la vejetanto. - EF.}
\l{}
\l{\sou{burleska}.\cel{2}Di qua la komikeso esas extravago. - DEFIS.}
\l{}
\l{\sou{burnuso}. Mantelo de lano blanka o bruna, kun kapuco quan}
\l{\cel{4}weras la Arabi. - DEFR.}
\l{}
\l{\sou{burokrato}. La persono qua molestoze uzas, a publiko, la povo}
\l{\cel{3}quan donas a lu lua grado en la kontori di stato-administr-}
\l{\cel{3}eyo. - DEFIS.}
\l{}
\l{\sou{burgo}. Saketo por portar moneto-peci en la posho. - DEFIS.}
\l{}
\l{\sou{buso}.\cel{2}(\sou{bot.}) Arboreto di la familio \textquotedbl{}euforbiacei\textquotedbl{}, di qua la}
\l{\cel{3}foliaro esas sempre verda. - L.\cel{1}\sou{buxus.}\cel{2}- DEFIRS.}
\l{}
\l{\sou{busho}. Grupo, asemblajo de arboreti (nefruktifera). - DEFI.}
\l{}
\l{\sou{bushelo}. En Francia, olima mezurilo di kapaceso (cirkume l3}
\l{\cel{3}litri) ; nun, la bushelo kontenas -segun la regioni) l o}
\l{\cel{3}2 dekalitri. - En Usa, la kapaceso varias segun la stati.}
\l{\cel{3}- EFR.}
\l{}
\l{\sou{busko.}\cel{2}Lamo flexebla ek barto-materio, stalo, e c., qua, per}
\l{\cel{3}aplikeso muldoza a la kurveso di la pektoro,\cel{1}\sou{man}tenas ri-}
\l{\cel{3}gide la parto avana di korsajo, di jupo-korpo. - EF.}
\l{}
\l{\sou{busolo}. Instrumento facita segun la duopla proprajo di la}
\l{\cel{3}magnet-agulo : direktar sua pinto ad (vers) la nordo-polo,}
\l{\cel{3}e proximigar su ad vertikalo segun ke lu portesas ad la}
\l{\cel{3}poli. - DFI.}
\l{}
\l{\sou{busprito}. (\sou{navig.)}\cel{2}Masto avana di navo, inklinita a la aquo-}
\l{\cel{3}surfaco, qua portas la seglo avana e fixigita an kordego,}
\l{\cel{3}- DEFIRS.}
\l{}
\l{\sou{busto}. I. Che la homo : la kapo e la parto supra di la korpo.}
\l{\cel{3}- II. Skulturo qua reprezentas la kapo e la parto supra}
\l{\cel{3}di la korpo, maxim-multa-kaze sen la brakii. - DEFIRS.}
\l{}
\l{\sou{bustrofedono}. Skribo-maniero che la Greki antiqua : la linei}
\l{\cel{3}de literi iras alterne de dextre a sinistre, e de sinistre}
\l{\cel{3}a dextre.\cel{2}Onu renkontras olu en ula epigrafi, subskriburi,}
\l{\cel{3}medalii. - DEF.}
\l{}
\l{\sou{butar}. (\sou{netrans., kontre)}.\cel{2}I. Marchante, pedo-shokar ne-expek-}
\l{\cel{3}tite kontre korpo qua salias sur la voyo. - II. Apogar ulo}
\l{\cel{3}kontre trabo, arko, e c. - EF.}
\l{}
\l{\sou{butiko}. Chambrego, ter-etaje, qua debushas an la strado, ed en}
\l{\cel{3}qua la komercisti expozas e vendas endetale sua vari. - FI.}
}
